%!TeX root=../pridetop.tex
\chapter[Chapter \thechapter]{}
	
	
\begin{figure}[t!]
\centering
\includegraphics[width=.8\linewidth]{45top}
\captionlistentry{Engaged by the river}
\end{figure}


\lettrine[lines=6,image=true]{initials/chap45c}{onvinced}  as Elizabeth now was that Miss Bingley's dislike of her had originated in jealousy, she could not help feeling how very unwelcome her appearance at Pemberley must be to her, and was curious to know with how much civility on that lady's side the acquaintance would now be renewed.

\zz
On reaching the house, they were shown through the hall into the saloon, whose northern aspect rendered it delightful for summer. Its windows, opening to the ground, admitted a most refreshing view of the high woody hills behind the house, and of the beautiful oaks and Spanish chestnuts which were scattered over the intermediate lawn.

In this room they were received by Miss Darcy, who was sitting there with Mrs Hurst and Miss Bingley, and the lady with whom she lived in London. Georgiana's reception of them was very civil, but attended with all that embarrassment which, though proceeding from shyness and the fear of doing wrong, would easily give to those who felt themselves inferior the belief of her being proud and reserved. Mrs Gardiner and her niece, however, did her justice, and pitied her.

By Mrs Hurst and Miss Bingley they were noticed only by a courtesy; and on their being seated, a pause, awkward as such pauses must always be, succeeded for a few moments. It was first broken by Mrs Annesley, a genteel, agreeable-looking woman, whose endeavour to introduce some kind of discourse proved her to be more truly well-bred than either of the others; and between her and Mrs Gardiner, with occasional help from Elizabeth, the conversation was carried on. Miss Darcy looked as if she wished for courage enough to join in it; and sometimes did venture a short sentence, when there was least danger of its being heard.

Elizabeth soon saw that she was herself closely watched by Miss Bingley, and that she could not speak a word, especially to Miss Darcy, without calling her attention. This observation would not have prevented her from trying to talk to the latter, had they not been seated at an inconvenient distance; but she was not sorry to be spared the necessity of saying much: her own thoughts were employing her. She expected every moment that some of the gentlemen would enter the room: she wished, she feared, that the master of the house might be amongst them; and whether she wished or feared it most, she could scarcely determine. After sitting in this manner a quarter of an hour, without hearing Miss Bingley's voice, Elizabeth was roused by receiving from her a cold inquiry after the health of her family. She answered with equal indifference and brevity, and the other said no more.

The next variation which their visit afforded was produced by the entrance of servants with cold meat, cake, and a variety of all the finest fruits in season; but this did not take place till after many a significant look and smile from Mrs Annesley to Miss Darcy had been given, to remind her of her post. There was now employment for the whole party; for though they could not all talk, they could all eat; and the beautiful pyramids of grapes, nectarines, and peaches, soon collected them round the table.

While thus engaged, Elizabeth had a fair opportunity of deciding whether she most feared or wished for the appearance of Mr Darcy, by the feelings which prevailed on his entering the room; and then, though but a moment before she had believed her wishes to predominate, she began to regret that he came.

He had been some time with Mr Gardiner, who, with two or three other gentlemen from the house, was engaged by the river; and had left him only on learning that the ladies of the family intended a visit to Georgiana that morning. No sooner did he appear, than Elizabeth wisely resolved to be perfectly easy and unembarrassed;—a resolution the more necessary to be made, but perhaps not the more easily kept, because she saw that the suspicions of the whole party were awakened against them, and that there was scarcely an eye which did not watch his behaviour when he first came into the room. In no countenance was attentive curiosity so strongly marked as in Miss Bingley's, in spite of the smiles which overspread her face whenever she spoke to one of its objects; for jealousy had not yet made her desperate, and her attentions to Mr Darcy were by no means over. Miss Darcy, on her brother's entrance, exerted herself much more to talk; and Elizabeth saw that he was anxious for his sister and herself to get acquainted, and forwarded, as much as possible, every attempt at conversation on either side. Miss Bingley saw all this likewise; and, in the imprudence of anger, took the first opportunity of saying, with sneering civility,—

»Pray, Miss Eliza, are not the ——shire militia removed from Meryton? They must be a great loss to \textit{your} family.«

In Darcy's presence she dared not mention Wickham's name: but Elizabeth instantly comprehended that he was uppermost in her thoughts; and the various recollections connected with him gave her a moment's distress; but, exerting herself vigorously to repel the ill-natured attack, she presently answered the question in a tolerably disengaged tone. While she spoke, an involuntary glance showed her Darcy with a heightened complexion, earnestly looking at her, and his sister overcome with confusion, and unable to lift up her eyes. Had Miss Bingley known what pain she was then giving her beloved friend, she undoubtedly would have refrained from the hint; but she had merely intended to discompose Elizabeth, by bringing forward the idea of a man to whom she believed her partial, to make her betray a sensibility which might injure her in Darcy's opinion, and, perhaps, to remind the latter of all the follies and absurdities by which some part of her family were connected with that corps. Not a syllable had ever reached her of Miss Darcy's meditated elopement. To no creature had it been revealed, where secrecy was possible, except to Elizabeth; and from all Bingley's connections her brother was particularly anxious to conceal it, from that very wish which Elizabeth had long ago attributed to him, of their becoming hereafter her own. He had certainly formed such a plan; and without meaning that it should affect his endeavour to separate him from Miss Bennet, it is probable that it might add something to his lively concern for the welfare of his friend.

Elizabeth's collected behaviour, however, soon quieted his emotion; and as Miss Bingley, vexed and disappointed, dared not approach nearer to Wickham, Georgiana also recovered in time, though not enough to be able to speak any more. Her brother, whose eye she feared to meet, scarcely recollected her interest in the affair; and the very circumstance which had been designed to turn his thoughts from Elizabeth, seemed to have fixed them on her more and more cheerfully.

Their visit did not continue long after the question and answer above mentioned; and while Mr Darcy was attending them to their carriage, Miss Bingley was venting her feelings in criticisms on Elizabeth's person, behaviour, and dress. But Georgiana would not join her. Her brother's recommendation was enough to insure her favour: his judgment could not err; and he had spoken in such terms of Elizabeth, as to leave Georgiana without the power of finding her otherwise than lovely and amiable. When Darcy returned to the saloon, Miss Bingley could not help repeating to him some part of what she had been saying to his sister.

»How very ill Eliza Bennet looks this morning, Mr Darcy,« she cried: »I never in my life saw anyone so much altered as she is since the winter. She is grown so brown and coarse! Louisa and I were agreeing that we should not have known her again.«

However little Mr Darcy might have liked such an address, he contented himself with coolly replying, that he perceived no other alteration than her being rather tanned,—no miraculous consequence of travelling in the summer.

»For my own part,« she rejoined, »I must confess that I never could see any beauty in her. Her face is too thin; her complexion has no brilliancy; and her features are not at all handsome. Her nose wants character; there is nothing marked in its lines. Her teeth are tolerable, but not out of the common way; and as for her eyes, which have sometimes been called so fine, I never could perceive anything extraordinary in them. They have a sharp, shrewish look, which I do not like at all; and in her air altogether, there is a self-sufficiency without fashion, which is intolerable.«

Persuaded as Miss Bingley was that Darcy admired Elizabeth, this was not the best method of recommending herself; but angry people are not always wise; and in seeing him at last look somewhat nettled, she had all the success she expected. He was resolutely silent, however; and, from a determination of making him speak, she continued,—

»I remember, when we first knew her in Hertfordshire, how amazed we all were to find that she was a reputed beauty; and I particularly recollect your saying one night, after they had been dining at Netherfield, »\textit{She} a beauty! I should as soon call her mother a wit.« But afterwards she seemed to improve on you, and I believe you thought her rather pretty at one time.«

»Yes,« replied Darcy, who could contain himself no longer, »but \textit{that} was only when I first knew her; for it is many months since I have considered her as one of the handsomest women of my acquaintance.«

He then went away, and Miss Bingley was left to all the satisfaction of having forced him to say what gave no one any pain but herself.

Mrs Gardiner and Elizabeth talked of all that had occurred during their visit, as they returned, except what had particularly interested them both. The looks and behaviour of everybody they had seen were discussed, except of the person who had mostly engaged their attention. They talked of his sister, his friends, his house, his fruit, of everything but himself; yet Elizabeth was longing to know what Mrs Gardiner thought of him, and Mrs Gardiner would have been highly gratified by her niece's beginning the subject.
